Our work sits at the meeting point of four literatures---unlearning algorithms, the
measurement of memorisation and forgetting, the verification of unlearning, and
game-theoretic statistics---and it is the gap between the middle two that
game-theoretic statistics turns out to close. We survey each in turn and then state
precisely how VOUCH differs from its nearest neighbours.

\subsection{Unlearning algorithms and their benchmarks}

Exact unlearning by architectural design, exemplified by sharded training
\citep{bourtoule2021machine}, retrains only the shards touched by a deletion; it is
sound but cannot be applied to a model that was not trained for it, which is the
situation every deployed language model is in. The practical response has therefore
been \emph{approximate} unlearning, and for large language models a compact toolbox
has emerged: gradient ascent on the forget set and its retain-regularised variant
GradDiff; negative preference optimisation \citep{zhang2024npo}, which replaces the
unbounded ascent objective with a bounded preference loss and which its authors
introduced precisely to avoid the catastrophic collapse that ascent invites; the
simplified reappraisal SimNPO \citep{fan2025simnpo}, which discards NPO's frozen
reference model in favour of a length-normalised margin objective and which we
include as a subject in Section~\ref{sec:exp-bench}; representation misdirection
\citep{li2024wmdp}, which acts on hidden activations rather than on the token loss and
which we add as a subject in Section~\ref{sec:exp-rmu}; and task-vector negation. These methods are evaluated on
purpose-built benchmarks---TOFU's fictitious-author question answering
\citep{maini2024tofu}, MUSE's six-way news-and-books evaluation \citep{shi2024muse},
and the hazardous-knowledge multiple-choice of WMDP \citep{li2024wmdp}---and we treat
all of them, methods and benchmarks alike, as \emph{subjects} that VOUCH certifies
rather than as competitors.

\subsection{Measuring memorisation and forgetting}

The canary methodology VOUCH depends on descends directly from the secret-sharer line
of work. \citet{carlini2019secret} insert random secret sequences into training data
and define \emph{exposure}, the canary's rank among all possible instantiations
ordered by model log-perplexity, converting memorisation into a calibrated number;
they show that memorisation is commonplace, that it is not explained by overfitting,
and that it is eliminated only by differentially private training. VOUCH borrows the
planted-canary device and changes one thing that turns out to matter for inference:
where the secret sharer plants a canary and \emph{measures} its exposure, we plant a
\emph{pair} and randomise which twin enters training, so that the measurement acquires
a null. Exposure answers ``how memorised is this string''; the ghost twin answers
``could this have happened by chance,'' which is the question a certificate must
answer and the one \citet{zhang2024mia} show cannot be answered without controlled
randomisation.

The strongest membership attacks in this lineage calibrate per example rather than
thresholding a raw score. \citet{carlini2022lira} show that comparing a target model's
confidence on an example against the confidences that shadow models trained with and
without it produce---a likelihood-ratio test they call LiRA---dominates score-threshold
attacks, particularly in the low false-positive regime that matters for a privacy
claim. LiRA is the natural adversary against which to stress-test a declared score
class, since it uses information no member of ours has: $N$ additionally trained
models. It is for the same reason not a candidate for inclusion in $\F$, because a
verifier held to $2Qm$ forward passes and no training cannot run it. We use it in
Section~\ref{sec:exp-lira} as an attack from outside the class, with a positive
control that establishes it is working before its null on a certified model is read as
meaning anything.

The complementary literature measures forgetting rather than memorisation.
\citet{jagielski2023measuring} quantify how the success of privacy attacks on a
training example decays as training continues past it, and find that standard models
do forget over time---but also prove that non-convex models can retain memorised data
indefinitely, and show that forgetting disappears entirely under deterministic
training. Their conclusion is the premise of ours: incidental forgetting is real but
uncontrolled, and therefore cannot substitute for a certificate.
\citet{barbulescu2024each} document, in the same spirit, how per-sequence memorisation
strength governs how hard a datum is to remove, which is the phenomenon our
repetition strata are designed to traverse.

\subsection{Evaluating unlearning: a literature of negative results}

What the unlearning literature has increasingly documented is that the metrics used
to declare success are unreliable, and this body of work is the direct motivation for
VOUCH. A first strand shows that apparent forgetting is often mere suppression.
\citet{lucki2024adversarial} recover ostensibly unlearned capabilities through
lightweight attacks; \citet{zhang2025quantization} show that quantising an unlearned
model to $4$ bits resurrects forgotten knowledge; \citet{deeb2024remove} recover most
pre-unlearning accuracy by fine-tuning on adjacent facts; and \citet{fan2025relearning}
formalise relearning attacks and connect robustness to sharpness-aware minimisation.
Closest to our P1 probe is \citet{hu2025jogging}, who show that fine-tuning an
unlearned model on a small, \emph{benign}, publicly available corpus that never
contains the forget targets is enough to jog the model's memory and restore both
hazardous knowledge and verbatim memorised passages. Their framing---obfuscation
versus removal---is precisely the distinction our relearn probe operationalises, and
their result is why we treat a post-relearning verdict as part of the certificate
rather than as an optional extra: on our benchmark tier at $\eps = 0.2$ the same probe
revokes NPO's certificate on one MUSE/GPT-2 seed and one MUSE/Phi-1.5 seed, and leaves
it intact on every TOFU cell (Section~\ref{sec:exp-bench}).

A second strand indicts the evaluation protocols themselves.
\citet{feng2025existing} demonstrate that current LLM-unlearning evaluations can be
made to both overstate and understate success. \citet{wang2025towards} show that
forget-quality metrics track surface output behaviour rather than retained knowledge
and collapse under red-teaming, and---more damagingly for cross-paper
comparison---that the reported unlearning/retention trade-off does not lie on a Pareto
frontier, so method rankings are artifacts of hyperparameter choice; they propose
recalibrating retained performance before comparing. \citet{yuan2025closer} make the
complementary point that ROUGE-style output matching hides degenerate behaviour, and
add metrics for token diversity, sentence semantics, and factual correctness. Both
diagnoses apply to VOUCH's subjects and shaped how we report them: we do not rank
unlearning methods against one another, we certify each one against a declared
tolerance and pair every verdict with a utility reading, precisely because a ranking
built on a non-Pareto trade-off would not survive a change of hyperparameters. Our own
worked example in Figure~\ref{fig:example} is an instance of the degeneracy
\citet{yuan2025closer} warn about---an NPO-unlearned model emitting
\texttt{LLLLLL\ldots} on canary prompts while its held-out loss is untouched---and it
is caught here not by a fluency metric but by the two-sided certificate, which treats
scoring conspicuously \emph{below} the ghost twin as leakage.

A third strand is mechanistic and, for our purposes, the most pointed.
\citet{wei2025forget} show that facts presumed forgotten persist through correlated
knowledge, and \citet{xu2026reversibility} distinguish genuine erasure from
suppression at the representation level, finding behaviour easily restored by minimal
fine-tuning. \citet{li2026beliefs} identify the mechanism they call the
\emph{squeezing effect}: gradient-ascent-style objectives push probability mass off
the target response but redistribute it onto nearby high-likelihood regions, typically
semantically equivalent rephrasings, so the content resurfaces under paraphrase and
commonly used metrics report success while semantically related information remains.
Their remedy couples the unlearning objective to the model's own high-confidence
generations. This result bears directly on the scope of what VOUCH certifies and we
want to be exact about the relationship. VOUCH is a \emph{verification} framework and
inherits the limits of its declared score class: our scores read the likelihood of the
literal secret span under a fixed set of query wrappers, so probability mass that has
migrated to a paraphrase is, by construction, mass our default class does not see. The
framework's response is structural rather than rhetorical---the class $\F$ is declared,
the certificate is explicitly a statement about $\F$, and any score a sceptic can name
can be added to $\F$ at the cost of a slower certificate and no loss of validity, since
Theorem~\ref{thm:cert} needs no multiplicity correction. Sections~\ref{sec:exp-outclass} and
\ref{sec:exp-lira} run attacks from \emph{outside} the declared class against
already-certified models and report where they exceed the tolerance. We also
implement the score their result calls for: Section~\ref{sec:exp-para} adds a
paraphrase-aware member to $\F$, which costs one to two per cent in pairs and changes
no verdict---though we are explicit there that our random-string secrets have no
semantic neighbourhood for mass to migrate into, so this tests the score's cost rather
than the squeezing effect itself, which would need canaries whose secret is
paraphrasable.

The recurring message across all of this is that \emph{evaluation}, not optimisation,
is the bottleneck---which is exactly the problem VOUCH addresses, and which motivates
our recoverability probes (R-VOUCH) and our insistence on pairing every certificate
with a utility reading.

\subsection{Verifying and certifying unlearning}

The earliest guarantees came from \emph{certified removal}
\citep{guo2020certified,sekhari2021remember,neel2021descent,ginart2019making}, which
bounds the distance between the unlearned and retrained models in a differential-
privacy-style $(\eps_u,\delta_u)$ sense. These results are elegant but confined to
convex or near-convex regimes, and their bounds become vacuous at the scale of modern
language models; we use them not as a baseline but as a consistency check, showing in
Proposition~\ref{prop:bridge} that any $(\eps_u,\delta_u)$-certified algorithm passes
VOUCH---a one-directional implication, as Remark~\ref{rem:onedirection} states.
The verification of \emph{arbitrary} unlearning has proven far harder.
\citet{thudi2022necessity} argue that approximate unlearning cannot be audited from
weights alone and that auditable definitions are a prerequisite for any guarantee;
\citet{zhang2024verification} show that a dishonest provider can defeat both existing
families of verification strategies while passing them; and, most consequentially for
our design, \citet{zhang2024mia} establish that membership-inference scores
\emph{cannot} prove that a model was trained on given data unless the inclusion of
that data was randomised---the soundness principle that our fair-coin canary design
is built to satisfy. Two further impossibility results bound what any behavioural
certificate can claim: \citet{yu2025impossibility} show that path-dependence makes
retrain equivalence unattainable for multi-stage pipelines, which is why VOUCH's null
is internally generated rather than referenced to a retrained model, and the
reversibility findings above imply that no behavioural method can certify
parameter-space erasure---a limit we state explicitly rather than paper over.

A handful of recent efforts move toward statistical certification, and these are our
closest neighbours. \citet{pandey2025gaussian} give a hypothesis-testing notion of
certified unlearning, but it is fixed-sample and tied to a Gaussian high-dimensional
regression regime rather than to LLM-scale behavioural auditing.
\citet{pawelczyk2025machine} use data-poisoning as a lens on unlearning failure but
provide no error-controlled certificate. Cryptographic proof-of-unlearning
\citep{eisenhofer2025verifiable} attests that the correct computation \emph{ran},
which is orthogonal to whether forgetting \emph{holds} statistically and with which
VOUCH composes---and which is the right tool against the special-casing provider our
threat model excludes. The recent survey of \citet{xue2026survey} organises this
fast-moving area and confirms that no existing method delivers what VOUCH does: an
anytime-valid, distribution-free certificate of bounded residual influence. The
broader programme is surveyed by \citet{liu2025rethinking}.

\subsection{Game-theoretic statistics}

The inferential engine we borrow is that of safe, anytime-valid inference: e-values,
test martingales, and betting confidence sequences
\citep{ramdas2023game,grunwald2024safe,shafer2021testing,waudbysmith2024,howard2021time}.
The defining property, traceable to Ville's inequality \citep{ville1939}, is that a
nonnegative martingale started at one exceeds $1/\alpha$ with probability at most
$\alpha$ \emph{at every stopping time simultaneously}; this is precisely what a
streaming, continuously-monitored deletion audit requires, and precisely what a
fixed-sample test cannot offer. The betting perspective
\citep{waudbysmith2024,orabona2016coin} additionally supplies confidence sequences
whose bets adapt online to near-optimal growth, and---the ingredient our corrected
certificate turns on---the without-replacement construction for finite-population
means, in which the null boundary is recentred at each step on the mean the unsampled
remainder would need to carry. That device is what lets Theorem~\ref{thm:cert} dispense
with independence across pairs entirely. The one adjacent application of these tools
is to privacy auditing: \citet{steinke2023privacy} audit differential
privacy in a single training run, and \citet{gonzalez2025sequential} and
\citet{csillag2026dpevalues} bring sequential e-value machinery to DP auditing. The
distinction is fundamental to the hypothesis structure. Privacy auditing
\emph{lower}-bounds the privacy loss of a \emph{training} algorithm by mounting a
successful attack; VOUCH \emph{upper}-bounds the residual influence \emph{after
unlearning} by showing that no attack in a declared class succeeds---an equivalence
test, not a significance test, run over ghost twins that have passed through the
unlearning pipeline. To our knowledge, VOUCH is the first application of
game-theoretic statistics to unlearning certification.

\subsection{Positioning}

The contrast with our nearest neighbours can be summarised sharply. Fixed-sample
paired or two-sample tests---including the KS-based forget-quality metric of
\citet{maini2024tofu} and the min-$k\%$ leakage metric of \citet{shi2024muse}---are
invalidated by optional stopping and typically require a retrained reference;
Section~\ref{sec:exp-descriptive} runs both on our own models and seeds and reports
where they and VOUCH disagree. The sound-membership results of \citet{zhang2024mia}
tell us \emph{why} controlled randomness is indispensable but stop at a fixed-$n$
membership proof rather than a sequential forgetting certificate; the secret-sharer
and forgetting-measurement lines \citep{carlini2019secret,jagielski2023measuring}
supply the canary device and the decay measurements but no null and no error control;
shadow-model evaluators \citep{hayes2024inexact} are descriptive and cost dozens of
fine-tunes; certified removal \citep{guo2020certified,sekhari2021remember} does not
scale; and DP auditing \citep{steinke2023privacy,gonzalez2025sequential} solves the
opposite, lower-bounding, problem. VOUCH is, to our knowledge, the first framework to
formulate unlearning certification as sequential equivalence testing and to deliver
exact, finite-sample, anytime-valid certificates of bounded residual influence at LLM
scale, with soundness inherited from a randomised paired-canary design rather than
from distributional assumptions or shadow models.
